\section{Nonzero-Gap Finite-Enclosure Obstructions}
\label{sec:nonzero-gap-finite-enclosure-optimization}
\label{sec:finite-enclosure}

The finite-enclosure method identifies a compact set that the C triangle
must contain, then proves that this set does not fit inside any open unit
equilateral triangle.

Put
\[
 h:=\frac{\sqrt3}{2}.
\]

The closed-disk notation and $\Lambda$ are defined in
Section~\ref{sec:introduction}. By \zcref{lem:compact-cover-margin}, a
compact set contained in an open unit triangle has enclosure side below
one. Conversely, $\Lambda(K)<1$ gives such an open unit enclosure by
slightly enlarging a smaller closed enclosure. Thus it is enough to
exclude every open unit candidate containing the same fixed set $K$.

\subsection{The two tasks in a finite-enclosure proof}
Fix the original seven triangles and call their arrangement $\mathcal U$.
A nonempty compact witness set $K=K(\mathcal U)$ is then chosen once and for all. There are
two separate tasks:
\[
 \underbrace{K\subset U_C}_{\text{forced containment}}
 \qquad\hbox{and}\qquad
 \underbrace{\Lambda(K)\ge1}_{\text{enclosure obstruction}}.
\]
The first makes $\Lambda(K)<1$ by the compact covering margin, contradicting
the second. Most witnesses are points missed by all open V triangles.
A structural anchor, such as the C midpoint, may instead be included because
it is already known to belong to $U_C$; such an anchor need not be missed
by every V triangle.

To prove the second task, suppose an \emph{arbitrary} open unit triangle
$U'$ contains the fixed $K$. Its orientation, offsets, and radial exits may
vary. The original reaches and the witness coordinates do not. In particular,
$U'$ need not cover the original boundary gaps or complete the original
skeleton cover. Any inequality obtained from original coverage must be
established before $U'$ is introduced and retained as a fixed input.

\begin{center}
\small
\renewcommand{\arraystretch}{1.15}
\begin{tabular}{@{}p{.27\linewidth}p{.65\linewidth}@{}}
\toprule
Object & What is fixed or required \\
\midrule
Original arrangement $\mathcal U$ & Its actual reaches and all coverage
hypotheses of the current case. \\
Witness set $K(\mathcal U)$ & Chosen from the original arrangement; never
recomputed from the candidate. \\
Candidate $U'$ & Required only to contain $K$; its own geometric parameters
are recomputed. \\
Replacement arrangement & A different operation, used later to preserve
skeleton coverage, not an enclosure candidate. \\
\bottomrule
\end{tabular}
\end{center}

The six-point construction uses actual gap endpoints and total radial
endpoints. The four-point construction uses a supported radial endpoint
and boundary anchors. The zero-gap construction instead uses uniform
capacity bounds to choose six radial points and three frontier points.
Their common feature is fixed containment, not a common coordinate formula.


\subsection{Neighboring capacities and fixed radial endpoints}
\label{subsec:new-neighbor-ray-capacity}

The following capacities answer a local question: after boundary anchors
have been prescribed, how far toward the center can a unit triangle reach?
They are maxima over a family of triangles, not the actual reaches of the
original triangle. A larger prescribed anchor shrinks the feasible family.
The neighboring capacities are needed because a point on one radial segment
may be covered by the V triangle at either adjacent vertex.

The neighboring-arm capacities are the variational quantities
\[
 C_+(a,b):=
 \max\left\{c\in[0,1]:
 \begin{array}{l}
 \exists T\ [T\text{ is a closed unit equilateral triangle},\\
 \{V_0,X_5(1-a),X_0(b),(1-c)V_1\}\subseteq T]
 \end{array}\right\},
\]
\[
 C_-(a,b):=C_+(b,a),
 \qquad
 0\le a,b\le1,
 \qquad
 a+b\le1.
\]
Only a bound, not the full piecewise neighboring capacity, is needed below.
We first compute the diagonal capacity.

\begin{lemma}[Sharp diagonal neighboring capacity]
\label{lem:direct-neighbor-diagonal}
For $0\le m\le1/2$, one has $C_+(m,m)=C_-(m,m)=1-m$.
\end{lemma}
\begin{proof}
Use oblique coordinates with metric $x^2+y^2-xy$. The four points are
$P=(0,0)$, $A=(m,0)$, $B=(0,m)$, and $Z=(c,1)$. The case $m=0$
is immediate. If $0<m\le1/2$ and $c>1-m$, their hull has cyclic order
$P,A,Z,B$. \zcref[S]{lem:support-cell-rotation} reduces the enclosure to its four outward
edge normals. At $PA$ and $BP$, the side lengths are $c+m$ and $1+m$,
both greater than one.

Embed $(x,y)$ as $(x-y/2,hy)$. At the other two edges take unnormalized
normals $n_A=(h,1/2-c+m)$ and $n_B=(h(m-1),m/2+c-1/2)$. Put
\[
\begin{aligned}
 D_A&=(c-m)^2-(c-m)+1,&N_A&=m+c^2-cm-c+1,\\
 D_B&=c^2+cm-c+m^2-2m+1,&N_B&=cm+c^2-c-m+1.
\end{aligned}
\]
Selecting the contained points $(A,Z,P)$ in the three normal directions
at $n_A$, and $(B,P,Z)$ at $n_B$, gives support lower bounds
$L_{AZ}\ge N_A/\sqrt{D_A}$ and $L_{ZB}\ge N_B/\sqrt{D_B}$.
No active-support selector is assumed. With $q=c+m-1>0$,
\[
\begin{aligned}
 N_A^2-D_A={}&q^2\big((q+1-3m)^2+4m^2-2m+1\big)\\
 &+q(1-2m)(6m^2-2m+1)+m^2(1-2m)^2>0,\\
 N_B^2-D_B={}&q^4+(2-2m)q^3+(m^2-4m+2)q^2\\
 &+(1-m)(1-2m)q>0.
\end{aligned}
\]
Here $4m^2-2m+1\ge3/4$, $6m^2-2m+1>0$, and
$m^2-4m+2\ge1/4$. The numerators are positive:
$N_A\ge(1+m)(3-m)/4$ and $N_B=1-m+(1-m)q+q^2$.
Thus every caliper requires side greater than one.
Equality is realized by the unit triangle
$\conv\{(-m,0),(1-m,0),(1-m,1)\}$. Reflection gives $C_-$.
\end{proof}

\begin{lemma}[Common-pair domination]
\label{lem:new-common-pair-domination}
\label{lem:appendix-common-pair-domination}
If $p,q\ge0$ and $p+q\le1$, then
\[
 C_+(p,q),C_-(p,q)\le1-\min\{p,q\}\le c_{\max}(p,q).
\]
\end{lemma}
\begin{proof}
Put $m=\min(p,q)$. Convexity shortens the boundary anchors to $(m,m)$,
so \zcref{lem:direct-neighbor-diagonal} gives $C_\pm(p,q)\le1-m$.
For the own-ray bound, use coordinates with origin $V_0$ and basis
$V_5-V_0,V_1-V_0$. If $q=m$, the unit triangle
$\conv\{(-m,0),(1-m,0),(1-m,1)\}$ contains the origin, $(p,0)$,
$(0,m)$, and the own-ray point $(1-m,1-m)$, since $p\le1-m$.
Thus $c_{\max}(p,q)\ge1-m$; reflect if $p=m$.
Increasing an anchor shrinks the feasible family, so capacities are antitone.
\end{proof}

\begin{lemma}[Strict own-ray bound with positive slack]
\label{lem:strict-own-ray-positive-slack}
If $a,b>0$ and $a+b<1$, then
\[
 c_{\max}(a,b)>1-\min(a,b).
\]
\end{lemma}
\begin{proof}
Put $m=\min(a,b)$, $M=\max(a,b)$, and $\delta=1-a-b>0$.
By \zcref{lem:new-common-pair-domination}, $c_{\max}(a,b)\ge1-m$.
The exact local formula \zcref{prop:new-exact-local-set} has two branches. In the quartic branch, the defining polynomial
$F(C,m)=C^4-C^2+mC-m^2$ satisfies
\[
 F(1-m,m)=-m\bigl((1-m)^3+m^2\bigr)<0.
\]
In the triangular branch, the polynomial
$H(C)=((1-\delta)^2-1)C^2+MC-M^2$ satisfies
\[
 H(1-m)=-\delta\bigl((1-m)(1-2m)+m(2-m)\delta\bigr)<0.
\]
Thus equality with $1-m$ is impossible on either selected capacity branch.
The non-strict bound proves the strict inequality. This argument does not
select a geometric component from a polynomial sign; it excludes equality
for the capacity already specified by the exact local theorem.
\end{proof}

\begin{definition}[Fixed total radial endpoint]
\label{def:fixed-total-radial-endpoint}
Parametrize $r_i$ by $Z_i(c)=(1-c)V_i$, measured from $V_i$ toward $O$.
Let $u_{j\to i}$ be the largest parameter in the actual positive trace of
$T_j$ on $r_i$ for $j=i-1,i+1$, with value zero when no such positive trace
exists. Put
\[
 \gamma_i:=\max\{C_i,u_{i-1\to i},u_{i+1\to i}\},\qquad
 d_i:=1-\gamma_i,\qquad \widehat P_i:=d_iV_i.
\]
\end{definition}

Thus $\widehat P_i$ is the innermost endpoint reached by any local V
triangle on $r_i$, rather than just the endpoint of the own V triangle.
This distinction is essential when a nonsupercritical V triangle supports
an adjacent radial segment.

\begin{lemma}[Total-endpoint radial forcing]
\label{lem:fixed-total-radial-forcing}
\label{lem:appendix-fixed-total-endpoint}
One has $0<d_i<1$, and every $sV_i$ with $0\le s\le d_i$ is missed
by all open V triangles. Under skeleton coverage these points belong to $U_C$.
\end{lemma}
\begin{proof}
No closed V triangle contains $O$: moving its interior distinguished vertex
slightly away from $O$ would contradict diameter one. Each own trace has
positive length, so $0<\gamma_i<1$ and $0<d_i<1$.
For $0<s\le d_i$, containment of $sV_i$ in a local open role would extend
its trace to an inward parameter strictly above $1-s\ge\gamma_i$.
This contradicts maximality, including the case of an absent positive
neighboring trace. Nonlocal roles are excluded by
\[
 \|sV_i-V_{i\pm2}\|^2=1+s+s^2>1,
 \qquad \|sV_i-V_{i+3}\|=1+s>1.
\]
The case $s=0$ is the already-excluded center. Skeleton coverage supplies $U_C$.
\end{proof}

\begin{corollary}[Clipped radial forcing]
\label{cor:clipped-radial-forcing}
For $j=i-1,i,i+1$, let $0\le a_j\le A_j$, $0\le b_j\le B_j$, and
$a_j+b_j\le1$. Define
\[
 s_i:=\min\{f(a_i,b_i),a_{i-1},b_{i-1},a_{i+1},b_{i+1}\}.
\]
Every $sV_i$ with $0\le s\le s_i$ is missed by all open V triangles.
For actual reaches on a gap-free nonsupercritical path, this simplifies to
\begin{equation}
 s_i=\min\{f(A_i,B_i),A_{i-1},B_{i+1}\}.
 \label{eq:clipped-path-radius}
\end{equation}
\end{corollary}
\begin{proof}
The own endpoint is at most $1-f(a_i,b_i)$. By common-pair domination,
each positive neighboring endpoint is at most $1-\min(a_j,b_j)$;
an absent trace contributes zero. Hence $\gamma_i\le1-s_i$ and
\zcref{lem:fixed-total-radial-forcing} applies. On a gap-free
nonsupercritical path, $A$ increases and $B$ decreases by the boundary
deficit identities. Together with $f(A_i,B_i)\le\min(A_i,B_i)$, these
inequalities remove $B_{i-1}$ and $A_{i+1}$ from the minimum.
\end{proof}

The actual radius $d_i$ comes from maximal traces; $s_i\le d_i$ is a
certified lower bound. Neither may be replaced by the own-ray radius
without accounting for neighboring suppliers.

\begin{corollary}[Uniform common-pair radial forcing]
\label{cor:new-uniform-common-pair-forcing}
Suppose $p,q\ge0$, $p+q\le1$, $p\le A_i$, $q\le B_i$ for every $i$,
and $0<c_*:=c_{\max}(p,q)<1$. Then $(1-c_*)V_i$ is missed by all
open V triangles, without a restriction on their normalized types.
Under skeleton coverage,
\[
 (1-c_*)V_i\in U_C\quad(0\le i\le5),
 \qquad \mathbb B(O,h(1-c_*))\subset U_C.
\]
\end{corollary}
\begin{proof}
Use the selected pair $(p,q)$ at every vertex in
\zcref{cor:clipped-radial-forcing}. Since $f(p,q)\le\min(p,q)$,
the clipped radius equals $1-c_*$. Convexity of $U_C$ supplies the
inner hexagon and its displayed inscribed disk.
\end{proof}

\subsection{Fixed witnesses and candidate enclosures}
\label{subsec:trace-exact-ab-envelopes}

Each witness set is fixed from the original V triangles before an
arbitrary enclosing candidate is chosen. The finite calipers vary their support orientation, but the witness points do not. Gap endpoints are supplied
by \zcref{lem:gap-exhaustion}; radial witnesses use the total endpoints,
including all permitted neighboring traces.







Replacement is a separate skeleton-preserving operation, not a fourth
witness construction. BC handles the center-aligned case. For an off-center
supercritical role, its midpoint supplier leads to D, replacement, or the
Vd2 perimeter obstruction.

% Active elementary proof from numbered source 2615.
\subsection{Elementary local capacity bounds}
Write $f(a,b)=1-c_{\max}(a,b)$ for $a,b\ge0$, $a+b\le1$, and put
$m=\min(a,b)$. Let $h=\sqrt3/2$ and let $c(m)\in[h,1]$ be the root of
\[
 F(C,m)=C^4-C^2+mC-m^2=0,\qquad \ell(m)=1-c(m).
\]
It is unique because $F_C>0$ on $[h,1]$,
$F(h,m)=-(m-h/2)^2\le0$, and $F(1,m)=m(1-m)\ge0$.
Define
\[
 q(m)=\frac{m(1-m)}2,\qquad
 \beta(m)=\frac{2m}{4+3m},\qquad \kappa(m)=1+\frac32m.
\]
\begin{lemma}[Own-ray deficit bounds]
\label{lem:bc-capacity-tools}
For $a,b\ge0$, $a+b\le1$, and $m=\min(a,b)$,
\[
 q(m)\le\ell(m)\le f(a,b)\le m.
\]
For $m\le1/4$, also $\beta(m)\le\ell(m)$ and $\beta(m)\ge q(m)$.
\end{lemma}
\begin{proof}
The upper bound and capacity antitonicity follow from
\zcref{lem:new-common-pair-domination}. By the selected local formula
\zcref{prop:new-exact-local-set}, for $m\le h/2$ the diagonal pair lies
in the quartic cell (its cell polynomial is $m^2(16m^2-3)\le0$), so
$c_{\max}(a,b)\le c_{\max}(m,m)=c(m)$. For $m\ge h/2$,
antitonicity gives $c_{\max}(a,b)\le c_{\max}(h/2,h/2)=h\le c(m)$.
Thus $f(a,b)\ge\ell(m)$.

For $c=c(m)$, the root equation and a factorization give
\[
 \ell c^2(1+c)=m(c-m),
\]
\[
 2(c-m)-(1-m)c^2(1+c)
 =(1-c)\bigl((1-m)c(c+2)-2m\bigr)\ge0.
\]
The bracket is at least
$\tfrac12\cdot\tfrac34\cdot\tfrac{11}4-1=1/32$, since $c\ge h>3/4$
and $m\le1/2$. Division by $c^2(1+c)>0$ proves $\ell\ge q(m)$.

For the stronger bound on $m\le1/4$, put
\[
 B(m)=81m^4+405m^3+656m^2+272m-128.
\]
Direct substitution gives
\[
 F(1-\beta(m),m)=-\frac{m^2B(m)}{(4+3m)^4}.
\]
The polynomial $B$ is increasing and $B(1/4)=-3163/256<0$.
Also $1-\beta(m)\ge17/19>h$. Monotonicity of $F$ in $C$ implies
$\beta(m)\le\ell(m)$. Finally
$\beta(m)-q(m)=m^2(3m+1)/(2(3m+4))\ge0$. At $m=0$ all bounds are immediate.
\end{proof}

\begin{lemma}[Quarter-range slack envelope]
\label{lem:bc-slack-envelope}
For $a,b\ge0$, $a+b\le1$, $m=\min(a,b)\le1/4$, and $\delta=1-a-b$,
\[
 f(a,b)\ge\max\{\beta(m),\ m-\kappa(m)\delta\}.
\]
The two entries meet at $\delta=\beta(m)$. This estimate is asserted only
on $m\le1/4$; the broader historical envelope is not needed here.
\end{lemma}
\begin{proof}
Since $m-\kappa(m)\beta(m)=\beta(m)$, the baseline in
\zcref{lem:bc-capacity-tools} proves the claim for $\delta\ge\beta(m)$.
The case $m=0$ is immediate. Suppose $m>0$ and $0<\delta\le\beta(m)$.
Put $M=1-m-\delta$. Then $M\ge m$ (indeed
$1-2m-\beta(m)\ge15/38>0$) and $1-\delta\ge1-\beta(m)\ge c(m)$.
The selected triangular branch of \zcref{prop:new-exact-local-set}
therefore applies: the capacity is the smaller root of
\[
 H(C)=\bigl((1-\delta)^2-1\bigr)C^2+MC-M^2.
\]
At $\overline C=1-m+\kappa(m)\delta$, expansion gives
$H(\overline C)=\delta N_m(\delta)/4$, where
\[
\begin{split}
 N_m(\delta)={}&(3m+2)^2\delta^3-10m(3m+2)\delta^2\\
 &+(28m^2-22m-20)\delta+14m(1-m).
\end{split}
\]
Because $\delta\le\beta(m)\le m/2\le1/8$,
\[
 N_m'(\delta)\le3(11/4)^2(1/8)^2+7/4-20=-18325/1024<0.
\]
At the right endpoint,
\[
 N_m(\beta(m))=-\frac{2mB(m)}{(4+3m)^3}>0.
\]
Thus $H(\overline C)>0$. The quadratic is concave and the capacity is
its selected smaller root, so $c_{\max}(M,m)<\overline C$. This proves
the affine lower bound. At $\delta=0$, the exact capacity is $1-m$.
\end{proof}

\begin{lemma}[Coupled nonuniform capacity inequality]
\label{lem:bc-coupled-capacity}
If $0<x\le y<1$ and $x/2\le z\le y$, put
\[
 r=f(1-y,z),\qquad t=f(1-z,x/2).
\]
Then
\[
 2(r+t)>(1-y+r)(2x-y+r).
\]
\end{lemma}
\begin{proof}
Put $u=x/2$ and
\[
 P(y,r,t)=2(r+t)-(1-y+r)(4u-y+r).
\]
The domain is $0<u$, $2u\le y<1$, $u\le z\le y$ and
$0<r,t\le1/2$. By \zcref{lem:bc-capacity-tools},
\[
 P_r=1-4u+2y-2r\ge1-2r\ge0,\qquad P_t=2.
\]
Hence valid lower bounds on either radius may be substituted. These are
scalar comparisons; the exact witness radii are not changed.

\readerheading{Easy regions.}
First let $y\ge1/2$ and put $a=1-y$. Then $0<a\le1/2$,
$u\le(1-a)/2$, and $u\le z\le1-a$.
If $z\ge a$ and $z\le1-u$, substitute $q(a),q(u)$.
The expression is concave in $u$; at $u=0$ it equals
$2q(a)+(a+q(a))(1-a-q(a))>0$, and at $u=(1-a)/2$ it equals
$(1-a)^3(1+a)/4>0$.
If $z\ge a$ and $z\ge1-u$, both radii are at least $q(a)$ because
$a\le1-z\le u\le1/2$ and $q$ is increasing. The substituted expression
decreases with $u$ and at $u=(1-a)/2$ equals
$a(1-a)(a^2-a+2)/4>0$.
If $z\le a$, both radii are at least $q(u)$. The expression is concave
in $a\in[u,\min(1/2,1-2u)]$, with endpoint values
\[
\begin{array}{c|c|c}
a&P\text{ after substitution}&u\text{ range}\\\hline
u&u(14-43u+14u^2-u^3)/4&(0,1/3]\\
1/2&(1-17u^2+10u^3-u^4)/4&(0,1/4]\\
1-2u&u(-u^3+2u^2+9u-2)/4&[1/4,1/3].
\end{array}
\]
The three brackets respectively decrease to $32/27$, decrease to
$23/256$, and increase from $23/64$, so all are positive.

Now let $y\le1/2$ and $z\ge2u$. Substitute $q(z),q(u)$.
The result increases with $y$, so set $y=z$; it is then concave in
$u\in[0,z/2]$. The value at zero is positive, and the value at
$u=z/2$ is $z^3(2-z)/4>0$.
Finally, for $y\le1/2$ and $1/4\le z\le2u$, the same substitution
allows $y=2u$, giving
\[
 P\ge q(z)(1-q(z))-u+3u^2
 \ge\frac3{32}\frac{29}{32}-\frac1{12}=\frac5{3072}>0.
\]
Here $q(z)(1-q(z))$ increases on $[1/4,1/2]$ and
$-u+3u^2=3(u-1/6)^2-1/12$.

\readerheading{The remaining region and its switches.}
It remains to treat
\[
 0<u\le1/4,\quad u\le z\le\min(2u,1/4),\quad 2u\le y\le1/2.
\]
The smaller capacity coordinates are $z$ and $u$. By
\zcref{lem:bc-slack-envelope} compare with
\[
 R_0=\max\{\beta(z),z-\kappa(z)(y-z)\},\qquad
 T_0=\max\{\beta(u),u-\kappa(u)(z-u)\}.
\]
Put $g(v)=v+\beta(v)$, an increasing rational function. The switches are
$y=g(z)$ and $z=g(u)$.
For fixed $u,z$, the affine branch of $R_0$ gives
\[
 \frac{d^2}{dy^2}P(y,R_0,T_0)=-2(1+\kappa(z))^2<0.
\]
The baseline branch gives $dP/dy=1+4u+2\beta(z)-2y>0$.
Since $g(z)\le g(1/4)=27/76<1/2$, only $y=2u$ or $y=g(z)$
needs to be checked, whenever the latter belongs to the domain.

On $y=2u$,
\[
 P=R_0(1-R_0)+2T_0-2u(1-2u).
\]
Split the $z$ interval at $g(z)=2u$ and $z=g(u)$.
The function $\beta(z)(1-\beta(z))$ is concave, since $\beta$ is concave
and $v(1-v)$ is increasing and concave on $[0,1/2]$.
On the other branch,
\[
 R=2z+\tfrac32z^2-2u-3uz,\qquad R'=2+3z-3u\ge2,\qquad R''=3.
\]
Where it is active, $0<R\le z\le1/4$, so
\[
 (R(1-R))''=3(1-2R)-2(R')^2\le-5<0.
\]
As $T_0$ is affine or constant on each piece, $P$ is piecewise concave.
At $z=u$ it equals $\beta(u)(1-\beta(u))+4u^2>0$; at $z=2u$ it
is $2T_0>0$. The endpoint $z=1/4$ is covered by the easy-region bound,
since the comparison radii dominate the $q$ bounds. Only the switches remain.

On $y=g(z)$,
\[
 P=2\beta(z)+2T_0-(1-z)(4u-z),\qquad z/2\le u\le g(z)/2.
\]
Its $u$ derivative on the baseline branch is
$16/(4+3u)^2-4(1-z)\le-2<0$; on the other branch it is $6u+z>0$.
Thus the minimum lies at an endpoint or $z=g(u)$. The endpoints are
$z=2u$, already checked, and transition A below. The ordering of switches
is not assumed; coincident switches are included.

\readerheading{Three transition checks.}
In A, $y=2u=g(z)$. Using the affine entry as a lower bound for $T_0$ gives
\[
 P\ge\frac{z^2(9z^2+36z+44)}{4(3z+4)^2}>0.
\]
In B, $y=2u$ and $z=g(u)$, so $T_0=\beta(u)$. If $R_B$ denotes the
affine entry of $R_0$, then
\[
 R_B-u(1-2u)=\frac{u^2(9u^2+6u+4)}{2(3u+4)^2}>0.
\]
Since $v(1-v)$ increases on $[0,1/2]$,
\[
 P\ge\frac{u^2(1+19u-4u^2-12u^3)}{3u+4}>0,
\]
where $4u^2+12u^3\le7/16$.
In C, $y=g(z)$ and $z=g(u)$, direct substitution gives
\[
 P=\frac{3u^2(81u^4+441u^3+768u^2+460u+48)}
 {(3u+2)(3u+4)^2(3u+8)}>0.
\]
These regions and transitions are exhaustive. Monotonicity in the radii
therefore proves the inequality for the exact capacities.
\end{proof}

% Geometry independent of the capacity construction.
\subsection{A five-point threshold criterion}
\begin{lemma}[Five-point threshold criterion]
\label{lem:bc-five-point}
\label{lem:five-point-threshold}
Suppose $0<x\le y<1$, $0<r\le\min(y,1-y)$, and $0<t\le x/2$.
Put $G_w=(1-w)V_0+wV_1$, $k=y-r$, and
$F=\{M_0,G_x,G_y,rV_2,tV_4\}$. Then
\begin{equation}
 \Lambda(F)>1\quad\Longleftrightarrow\quad
 r+t+\max\{1/2,1-x(1-k)\}>\sqrt{1-k+k^2}.
 \label{eq:five-point-threshold}
\end{equation}
The gap points may coincide. This is a criterion at side one, not a
claim that the displayed caliper always minimizes the enclosing side.
\end{lemma}
\begin{proof}
Put $M=M_0$, $P=rV_2$, and $Q=tV_4$. For $x<y$, these points have
hull order $M,G_x,G_y,P,Q$. The stated bounds make the consecutive
cross products positive. At $G_yP$, the outward normal $n=(-hk,1-k/2)$ has
squared norm $1-k+k^2$. The three supports divided by $h$ are
$r$, $t$, and $\max\{1/2,1-x(1-k)\}$. Hence its caliper side is
\[
 L_3=\frac{r+t+\max\{1/2,1-x(1-k)\}}{\sqrt{1-k+k^2}}.
\]
The other four calipers exceed one. At the edges $MG_x$, $G_xG_y$,
and $QM$, selected point projections give
\[
 L_1>\frac{1-x+2x^2}{\sqrt{1-2x+4x^2}},\qquad
 L_2=1+r+t,\qquad
 L_5\ge\frac{1+t+2t^2}{\sqrt{1+2t+4t^2}}.
\]
For $L_1$ use $y\ge x$ and the strictly positive contribution $2tx$;
for $L_5$ use $y\ge x\ge2t$. The squared differences are
\[
 x^2(2x-1)^2\ge0,\qquad t^2(1+2t)^2>0,
\]
respectively. At $PQ$, put $v=t/r>0$. Selected projections give
\[
 L_4\ge\frac{t+v+x+\max\{1/2,1-(1+v)x\}}{\sqrt{1+v+v^2}}.
\]
The piecewise-linear expression in $x$ is minimized at $1/[2(1+v)]$.
Thus its numerator exceeds $g(v)=v+1/2+1/[2(1+v)]$, and
\[
 g(v)^2-(1+v+v^2)=\frac{v^2}{4(1+v)^2}>0.
\]
By \zcref{thm:cert-caliper}, only $L_3$ can cross the unit threshold.
When $x=y$, discard the zero edge $G_xG_y$; the other four calculations
remain valid, without a limiting strictness argument.
\end{proof}

\readerheading{Capacity instance.}
For $x/2\le z\le y$, take $r=f(1-y,z)$ and $t=f(1-z,x/2)$.
The bounds of \zcref{lem:bc-capacity-tools} meet the geometric hypotheses,
and \zcref{lem:bc-coupled-capacity} gives $r+t>(1-k)(x-k/2)$.
Consequently the left side of \eqref{eq:five-point-threshold} exceeds
$1-k/2+k^2/2\ge\sqrt{1-k+k^2}$; the squared difference is
$k^2(1-k)^2/4$. Thus this capacity instance has $\Lambda(F)>1$.


\subsection{A selected gap forces five points (BC)}

Normalize the center midpoint to $M_0$ and choose any actual gap, reflected
to $J_0$. The other incident edge $e_{5,0}$ may or may not contain a gap.
Only the four middle edges $e_{1,2},\ldots,e_{4,5}$ must be gap-free.
The following transfer supplies the tail bound needed by the path argument.

\begin{lemma}[Shared-gap-anchor boundary transfer]
\label{lem:shared-gap-anchor-transfer}
If the original open roles cover $\partial H$ and $J_0$ is an actual gap,
then, independently of the gap status of $e_{5,0}$,
\[
 B_5>1-M_0(B_0)>B_0/2.
\]
Here $M_0(z)$ is the scalar diameter envelope.
\end{lemma}
\begin{proof}
Put $G=X_0(B_0)$ and $L=X_5(B_5)$. Actual gap forcing gives
$G\in U_C$, while maximality gives $G\in T_0$ and $L\notin U_5$.
Boundary locality and the original perimeter cover give
$L\in U_0\cup U_C$. Thus one unit triangle contains $L$ openly and
$G$ in its closure, so $\|G-L\|<1$. Indeed, at equality one could move
$L$ slightly away from $G$ inside that open triangle and exceed its diameter.
The two incident-edge directions at $V_0$ meet at $120$ degrees, hence
\[
 B_0^2+B_0(1-B_5)+(1-B_5)^2<1.
\]
Apply \zcref{lem:signed-diameter-transfer}; its last inequality is strict
because $B_0>0$. The proof includes singleton gaps and does not require
$L$ to be a witness in the final finite set.
\end{proof}

In the pure enclosure statement below the tail bound $B_5\ge B_0/2$
is explicit. In the covering application, the preceding lemma establishes
it for the original roles before an arbitrary enclosing candidate is chosen.
We do not assume that the candidate also covers the discarded gap.

Define the two nonuniform comparison pairs and radii by
\[
 x=B_0,\qquad y=1-A_1,\qquad z=B_3.
\]
\[
 r=1-c_{\max}(1-y,z),\qquad t=1-c_{\max}(1-z,x/2).
\]
Put $\widehat t=\min\{t,A_3\}$ and fix
\[
 K_{BC}:=\{M_0,X_0(x),X_0(y),rV_2,\widehat tV_4\}.
\]
There are at most five points, including singleton gaps. The midpoint is a
structural anchor; the other points are missed by every open V role.
The cap by $A_3$
controls the remaining neighboring supplier on $r_4$; it is not silently dropped.

\readerheading{The five-point selected-gap construction.}
Set $x=B_0$, $y=1-A_1$, and $z=B_3$. The two boundary pairs
$(1-y,z)$ and $(1-z,x/2)$ are different; they preserve the shared boundary
parameter $z$ instead of replacing all path data by a uniform minimum.
Their capacity deficits give $r$ and $t$. The five witnesses are
$M_0,X_0(x),X_0(y),rV_2,\min(t,A_3)V_4$.
There is no point on $r_3$, although $B_3$ remains part of the construction.
The cap by $A_3$ handles a neighboring supplier. When it is active,
a pair of witnesses already has distance greater than one; otherwise the capacity inequality verifies the single decisive caliper
of the five-point threshold criterion. Singleton gaps merely identify
$X_0(x)$ and $X_0(y)$.


\begin{theorem}[Unified five-point selected-gap enclosure]
\label{prop:new-nplus-one-all-vd0}
\label{lem:appendix-fixed-transverse}
Suppose $J_0$ is an actual gap, $A_i+B_i\le1$ for $1\le i\le5$, and
\[
 B_i+A_{i+1}>1\quad(1\le i\le4),\qquad B_5\ge B_0/2.
\]
Then $\Lambda(K_{BC})>1$. Under skeleton coverage with distinguished
C midpoint $M_0$, the fixed set $K_{BC}$ lies in $U_C$.
There is no condition on the gap status of $e_{5,0}$ or supercriticality
of $T_0$, and no type restriction on the nonsupercritical path.
\end{theorem}
\begin{proof}
The boundary-deficit identities give increasing $A_1,\ldots,A_5$ and
decreasing $B_1,\ldots,B_5$. Thus $0<x\le y<1$ and $x/2\le z\le y$.
Put $\rho_i=f(A_i,B_i)$. By \zcref{cor:clipped-radial-forcing},
the safe radii are $\min\{\rho_i,A_{i-1},B_{i+1}\}$ for $i=2,3,4$.
The bounds $(A_2,B_2)\ge(1-y,z)$ imply
$r\le\rho_2$, $r\le A_1$, and $r\le B_3$. Hence $rV_2$ is safe.
Similarly $(A_4,B_4)\ge(1-z,x/2)$ gives $t\le\rho_4$, and
$t\le x/2\le B_5$; clipping by $A_3$ makes $\widehat tV_4$ safe.
The gap endpoints follow from \zcref{lem:gap-exhaustion}, and the midpoint
belongs to $U_C$ by hypothesis. All witnesses are fixed before enclosure is tested.

If $A_3\ge t$, apply the capacity instance of \zcref{lem:bc-five-point}. Otherwise $A_3<t$ and
$A_3\ge A_1=1-y$, so
\[
 \|X_0(y)-A_3V_4\|^2\ge
 \|X_0(y)-(1-y)V_4\|^2=1+(1-y)(2-y)>1.
\]
The diameter obstruction gives $\Lambda(K_{BC})>1$ in this case also.
No signed center coordinates or CE1 return are used in this enclosure argument.
\end{proof}

\begin{theorem}[Center-aligned path obstruction]
\label{thm:center-aligned-path}
Suppose a skeleton cover has at least one actual gap and its C triangle's
unique midpoint is $M_k$. If every $T_i$ with $i\ne k$ is nonsupercritical,
the configuration is impossible.
\end{theorem}
\begin{proof}
Normalize $k=0$. The signed normal form confines possible gap edges to
$e_{5,0},e_{0,1}$. Select either actual gap and reflect it to $J_0$.
The four middle edges are gap-free, and
\zcref{lem:shared-gap-anchor-transfer} gives $B_5>B_0/2$.
Apply \zcref{prop:new-nplus-one-all-vd0}. The midpoint belongs openly
to $U_C$, and all other witnesses are center-forced. Compact-open
containment contradicts $\Lambda(K_{BC})\ge1$ in either gap rank.
\end{proof}

\begin{theorem}[Every skeleton cover has a supercritical V triangle]
\label{prop:new-nplus-zero-gap-closures}
\label{thm:fixed-nzero}
A skeleton cover by the seven original open roles satisfies $N_+\ge1$.
\end{theorem}
\begin{proof}
If $N_+=0$ and there are no gaps, strict overlaps give
$6<\sum_i(A_i+B_i)\le6$. If a gap is present, all six roles are
nonsupercritical, so \zcref{thm:center-aligned-path} applies. This proof uses only BC; replacement will invoke N0, not conversely.
\end{proof}

\begin{lemma}[A unique midpoint supplier]
\label{lem:midpoint-supplier}
In a nonzero-gap skeleton cover, the preceding N0 theorem and the skeleton
budget leave only $N_+=1$, $N_{\rm sp}\le1$. Write $M_k$ for the C
midpoint and $\sigma$ for the supercritical index. If $\sigma\ne k$,
there is a unique positive-support role $T_\tau$, where
\[
 \tau\in\{\sigma-1,\sigma+1\},\qquad M_\sigma\in U_\tau.
\]
If $T_\tau$ is T3-like, then $\tau=k$. If $\tau\ne k$, the edge shared
by $T_\tau,T_\sigma$ is center-free, in both CE1 and CE2.
\end{lemma}
\begin{proof}
N0 excludes $N_+=0$. Two supercritical roles would require a distinct
positive-support rescuer by \zcref{lem:positive-support-rescuer}, and
\zcref{thm:common-skeleton-count} would exclude them. The same count
bound gives $N_{\rm sp}\le1$.

If $\sigma\ne k$, neither the C triangle nor $T_\sigma$ contains
$M_\sigma$. Diameter excludes nonlocal roles; a Vd0 role has no positive
adjacent trace and cannot contain that point openly. Hence its unique
supplier is an adjacent positive-support role. A T3-like supplier also
misses its own midpoint by \zcref{lem:t3-nonsupercritical}. Unless $\tau=k$, neither the center nor the
other roles could cover this second midpoint: all other roles are Vd0.
Finally, when $k\notin\{\tau,\sigma\}$, the shared edge is not incident
to $V_k$, whereas every open center boundary trace is incident to $V_k$
by the signed normal form.
\end{proof}

\begin{samepage}
\subsection{A supported midpoint forces four points (D)}

Let $Y(t)=(1-t)V_0+tV_5=X_5(1-t)$. The geometric lemma below is stated
independently of any V type or covering configuration.

\begin{theorem}[Four-point rescuer geometry]
\label{thm:fixed-four-point-rescuer}
\label{lem:appendix-fixed-four-point}
Suppose $a\ge0$, $\varepsilon>0$, $\beta\ge0$, and
\[
 s=a+\varepsilon,\qquad \beta\le\frac{\varepsilon}{s}.
\]
Then
\[
 \Lambda\{O,\varepsilon V_1,Y(a),Y(1-\beta)\}\ge1.
\]
Neither $a\le\varepsilon$ nor $s\le1$ is required by this geometric statement.
\end{theorem}
\begin{proof}
If $a=0$, the set contains $O,V_0$, at distance one. If $s\ge1$, then
\[
 \|Y(a)-\varepsilon V_1\|^2=1-s+s^2\ge1.
\]
These are diameter obstructions. Otherwise put $v=a/s$, so $0<s,v<1$.
The ratio gives $1-\beta\ge v$, while $a=sv\le v$. Hence $Y(v)$ lies
on $[Y(a),Y(1-\beta)]$. By convexity it suffices to enclose the quadrilateral,
in cyclic order,
\[
 O,\qquad Q=Y(v),\qquad A=Y(sv),\qquad P=s(1-v)V_1.
\]
With $h=\sqrt3/2$, its outward edge normals may be chosen as
\[
 (-hv,v/2-1),\quad v(1-s)(h,-1/2),\quad
 (hs,1-s/2),\quad s(1-v)(-h,1/2).
\]
The maximizing triples in each direction and its two $120$-degree rotations
are $(O,A,P)$, $(A,P,O)$, $(A,O,Q)$, and $(O,Q,A)$, respectively.
The resulting caliper sides are
\[
 \frac1{\sqrt{1-v+v^2}},\qquad 1+s(1-v),\qquad
 \frac1{\sqrt{1-s+s^2}},\qquad 1+v(1-s).
\]
All exceed one for $0<s,v<1$. The finite-caliper theorem
\zcref{thm:cert-caliper} completes the nondegenerate case.
\end{proof}

\readerheading{Why the ratio conditions matter.}
The identity in the preceding proof combines the supported radial point
with one boundary point to produce a point on $r_0$. Since $a\le\varepsilon$,
its distance from $O$ is $\varepsilon/s\ge1/2$. A candidate also
contains $O$, so it must contain $M_0$. The other boundary point then demands
more of the candidate's left boundary trace than its signed-center geometry
allows. This is why the application needs a ratio estimate: it converts the
local supplier geometry into this same four-point contradiction.


\end{samepage}

\begin{lemma}[Scalar rescuer ratio test]
\label{lem:rescuer-ratio-test}
For $0\le x,c\le1/2$ and $M=(c+\sqrt{c^2-8c+4})/2$,
\[
 x\le1-M\quad\Longleftrightarrow\quad x^2+(c-2)x+c\ge0.
\]
\end{lemma}
\begin{proof}
The smaller root is $1-M$, while the other root is at least one.
The quadratic is nonnegative on $[0,1/2]$ exactly before its smaller root.
\end{proof}

\subsection{The supported-rescuer adapter}



\begin{lemma}[Supported endpoint to four-point enclosure]
\label{lem:new-rescuer-tail-budget}
Suppose $T_0$ has actual boundary endpoint $a=A_0$ toward $V_5$ and
supported interval $[c,u]$ on $r_1$, measured from $V_1$ toward $O$.
Put $\varepsilon=1-u>0$ and $M=M_c^{\rm sup}$. Assume
\[
 \varepsilon V_1\in U_C,\qquad A_0+\varepsilon\le1,
 \qquad \frac{A_0}{A_0+\varepsilon}\le1-M.
\]
If $T_1$ is uniquely supercritical with $C_1\ge c$, and the
four center-free edges from $e_{1,2}$ through $e_{4,5}$ join nonsupercritical
$T_2,T_3,T_4,T_5$, then the skeleton is not covered. The conclusion applies
to both possible nonzero gap ranks.
\end{lemma}
\begin{proof}
The strict-supercritical envelope gives $B_1<M$, and the four handoffs give
$B_5\le B_4\le B_3\le B_2\le B_1$. Thus $B_5<1-A_0$, so the left-edge
actual gap has endpoints $Y(A_0),Y(1-B_5)$, with $Y(t)=(1-t)V_0+tV_5$.
Under coverage these endpoints and $O,\varepsilon V_1$ belong to $U_C$.
Since $M\ge1/2$, the hypotheses imply
\[
 A_0\le\varepsilon,\qquad
 B_5<M\le\frac{\varepsilon}{A_0+\varepsilon}.
\]
Apply \zcref{thm:fixed-four-point-rescuer} with geometric parameters
$a=A_0$, $\beta=B_5$. The fixed set has enclosure number at least one,
contrary to compact containment in $U_C$.
\end{proof}

\subsection{T3-like supported-tail adapter}

\begin{lemma}[Solved T3-like supported-tail adapter]
\label{lem:appendix-t3-supported-tail}
Assume a nonzero-gap skeleton cover has center midpoint $M_0$, a T3-like
role $T_0$ supplying $M_1$, uniquely supercritical $T_1$, and nonsupercritical
Vd0 roles at the other four vertices. The actual supported interval
satisfies the hypotheses of \zcref{lem:new-rescuer-tail-budget}.
\end{lemma}
\begin{proof}
Use the original Type-II chart \eqref{eq:orientation-type-II}, with
positive first slack $\alpha_T$ and actual reach $a=A_0$.
Thus $0<t<1$, $z=\sqrt{1-t+t^2}$, and substitution on $r_1$ gives
\begin{equation}
 c=\frac{1+\alpha_T+a-z}{1-t},\qquad
 u=a+t,\qquad \varepsilon=1-u.
 \label{eq:t3-actual-supported-interval}
\end{equation}
Open midpoint containment gives $0<c<1/2<u<1$.
The supercritical own role ends before the midpoint and the other roles
have no positive trace on $r_1$. Thus $u$ is the total radial frontier:
$\varepsilon V_1\in U_C$ by \zcref{lem:fixed-total-radial-forcing}.
Coverage of the near endpoint gives $C_1\ge c$, since the center exits
before $M_1$.

Put $x=a/(1-t)$ and $\theta=t/(1+z)$. Then
\begin{equation}
 a+\varepsilon=1-t<1,\qquad
 c=x+\theta+\frac{\alpha_T}{1-t}>x+\theta,\qquad
 \frac a{a+\varepsilon}=x.
 \label{eq:t3-actual-ratio}
\end{equation}
In particular $0<x<c<1/2$. Since $z<1$ and
$x+(1-x)t=a+t>1/2$,
\[
 \theta>\frac t2>\frac{1-2x}{4(1-x)}.
\]
The quadratic $Q_c(x)=x^2+(c-2)x+c$ increases with $c$, so
\[
 Q_c(x)>2x^2+(\theta-1)x+\theta
 >\frac{(1-2x)(4x^2-3x+1)}{4(1-x)}>0.
\]
Here $4x^2-3x+1=4(x-3/8)^2+7/16$.
The scalar rescuer test \zcref{lem:rescuer-ratio-test} gives $x\le1-M_c^{\rm sup}$.
Together with the actual size relation and the forced endpoint, these
are exactly the common D hypotheses. No translated near endpoint is used.
\end{proof}

\begin{lemma}[Solved Vd1 supported-tail adapter]
\label{lem:appendix-vd1-supported-tail}
In the center-based Vd1 supplier placement, with $M_1\in U_0$ and
$T_1$ uniquely supercritical, the actual corner data imply the size and
ratio inequalities required by \zcref{lem:new-rescuer-tail-budget}.
\end{lemma}
\begin{proof}
Set $a=A_0$, $b=B_0$. By \zcref{lem:vd1-forward-orientation}, the
original corner parameter satisfies $t\ge1$. Put
$d=\sqrt{t^2+t+1}$ and write the actual supported interval as $[c,u]$.
Then $0<c<1/2<u<1$ and
\[
 c=\frac{t(1-b)}{t+1},\qquad
 u=\frac{d-a-tb-1}{t}.
\]
The midpoint condition and $tb=t-c(t+1)$ imply
\[
 a\le F(t,c):=d-1-\frac{3t}{2}+c(t+1).
\]
Since $d'<1$ and $c<1/2$, $F$ decreases with $t$, so
\[
 a\le F(t,c)\le F(1,c)=\sqrt3-\frac52+2c=:L(c).
\]
Put $\varepsilon=1-u$. Direct substitution gives
\[
 \varepsilon=\varepsilon_0+\frac at,\qquad
 \varepsilon_0=\frac{2t+1-c(t+1)-d}{t}>0.
\]
For the last sign, $d<t+1$ gives
$t\varepsilon_0>(1-c)t-c\ge1-2c>0$.
The function $v/(v+\varepsilon_0+v/t)$ increases for $v\ge0$, and
$\varepsilon_0+F(t,c)/t=1/2$. Hence, with $x=a/(a+\varepsilon)$,
\[
 x\le\frac{F(t,c)}{F(t,c)+1/2}\le2F(t,c)\le2L(c)
 <4c-\frac32,
\]
using $\sqrt3<7/4$. Thus $0<x<1/2$ and $c>(3+2x)/8$. Therefore
\[
 x^2+(c-2)x+c>
 x^2+\left(\frac{3+2x}{8}-2\right)x+\frac{3+2x}{8}
 =\frac{(1-2x)(3-5x)}8>0.
\]
The scalar rescuer test \zcref{lem:rescuer-ratio-test} proves $x\le1-M_c^{\rm sup}$.

Positive forward support also gives
\[
 (t+1)a+tb<(t+1)(d-1)-t^2<d-1,
\]
because $d<t+1$. Thus $u>a$ and $a+\varepsilon=1+a-u<1$.
As in the T3-like placement, $u$ is the total frontier and the near
endpoint is supplied by $T_1$, so $\varepsilon V_1\in U_C$ and
$C_1\ge c$. These are all the common D hypotheses.
\end{proof}

For the covering application, let $T_0$ supply $M_1$, with supported
interval $[c,u]$ on $r_1$, measured from $V_1$ toward $O$, and put
$\varepsilon=1-u$ and $P_T=\varepsilon V_1$. All interval endpoints here
belong to the original triangle, not to its translated majorant.
The local T3-like and Vd1 calculations establish
\[
 A_0+\varepsilon\le1,\qquad
 \frac{A_0}{A_0+\varepsilon}\le1-M_c^{\rm sup}.
\]
Together with $C_1\ge c$, the boundary path, and the forced endpoint,
\zcref{lem:new-rescuer-tail-budget} gives the single fixed witness
\[
 K_D=\{O,\varepsilon V_1,Y(A_0),Y(1-B_5)\}\subset U_C,
 \qquad \Lambda(K_D)\ge1.
\]
The two local verifications are
\zcref{lem:appendix-t3-supported-tail,lem:appendix-vd1-supported-tail};
the enclosure step does not distinguish one from two gaps.



\begin{figure}[!htbp]
\centering
\begin{minipage}[t]{.47\linewidth}
\centering
\resizebox{.96\linewidth}{!}{%
\begin{tikzpicture}[x=6cm,y=1cm,>=Latex]
  \node[panellabel,anchor=west] at (-.02,2.30) {(a) the supported interval on $r_1$};
  \draw[flow] (0,1.15)--(1,1.15);
  \fill (0,1.15) circle (.017) node[left=2pt,figlabel] {$V_1$};
  \fill (1,1.15) circle (.017) node[right=2pt,figlabel] {$O$};
  \coordinate (C) at (.24,1.15);
  \coordinate (M) at (.50,1.15);
  \coordinate (U) at (.73,1.15);
  \draw[actualreach] (C)--(U);
  \draw[VertexOrange,fill=white,line width=.7pt] (U) circle (.025);
  \fill[DemandRed] (U) circle (.011);
  \fill[black] (C) circle (.019) node[above=3pt,figlabel] {$c$};
  \fill[CenterBlue] (M) rectangle +(0.026,0.026);
  \node[above=3pt,figlabel] at (M) {$M_1$};
  \node[above=3pt,figlabel] at (U) {$u$};
  \draw[dimline] ($(U)+(0,-.22)$)--($(1,1.15)+(0,-.22)$)
    node[midway,below=1pt,figlabel] {$1-u$};
  \node[figlabel,align=center] at (.52,.12)
    {$P_T=(1-u)V_1$ is the O-side endpoint\\
     coordinate on the top axis is measured from $V_1$ toward $O$};
\end{tikzpicture}%
}
\end{minipage}\hfill
\begin{minipage}[t]{.47\linewidth}
\centering
\resizebox{.96\linewidth}{!}{%
\begin{tikzpicture}[scale=2.15,>=Latex]
  \node[panellabel,anchor=west] at (-1.08,1.22) {(b) why $P_T$ is center-forced};
  \HexCoords
  \DrawHexagon
  \DrawRays
  \coordinate (PT) at ($(O)!.27!(V1)$);
  \coordinate (M1) at ($(O)!.5!(V1)$);
  \draw[vertexrole]
    (V0)--($(V0)!.76!(V1)$)--($(O)!.20!(V1)$)--cycle;
  \draw[DemandRed,dashed,line width=.9pt]
    (V1)--($(V1)!.48!(V0)$)--($(O)!.48!(V1)$)--cycle;
  \draw[actualreach] ($(V1)!.22!(O)$)--($(V1)!.73!(O)$);
  \fill[DemandRed] (PT) circle (.025) node[left=2pt,figlabel] {$P_T$};
  \fill[CenterBlue] (M1) rectangle +(0.025,0.025);
  \node[figlabel,right] at (M1) {$M_1$};
  \draw[flow] (PT)--($(PT)!1.55!(O)$);
  \node[figlabel,align=left,anchor=west] at (-.96,-.72)
    {$T_0$: open endpoint at $P_T$\\
     $T_1$: misses $M_1$, hence the O-side segment\\
     other V roles: excluded by type or diameter};
  \node[figlabel,atlasbox,align=center] at (.48,-1.05)
    {$P_T\in U_C$};
\end{tikzpicture}%
}
\end{minipage}
\caption{Construction of the T3-like rescuer endpoint.  The interval
$[c,u]$ is parametrized from $V_1$ toward $O$, so its O-side endpoint is
$P_T=(1-u)V_1$.  The T3-like open role omits its endpoint, the adjacent
supercritical role contains $V_1$ but misses $M_1$ and therefore misses the
O-side segment by convexity, and all remaining V roles are excluded.  Thus
$P_T$ is forced into the C triangle.}
\label{fig:t3-rescuer-endpoint}
\end{figure}

\FloatBarrier

\readerheading{A change of configuration, not a new candidate.}
The away-supplier case uses a different operation from finite enclosure.
Under the stated two-chart hypotheses, two V triangles are replaced while
the distinguished roles and coverage of the skeleton are retained. The
result in this unique-supercritical case has no supercritical V triangle, contradicting
\zcref{thm:fixed-nzero}. This explains why that theorem was proved for
\emph{skeleton} covers before replacement was introduced.

Neither coverage of all of $H$ nor the original number of gaps is claimed
to survive. The contradiction uses only the preserved skeleton cover.
Likewise, the earlier translated closed-trace majorant is not itself a
valid replacement open triangle: its distinguished vertex lies on its
boundary. These three objects---a majorant, an enclosing candidate, and a
replacement arrangement---have different uses.


\subsection{Vd1 two-chart replacement}

The replacement below uses two different vertex charts.  It preserves the
five skeleton pieces that can change, then invokes the all-nonsupercritical
obstruction without assuming preservation of the input gap rank.

\begin{lemma}[Half-square admissibility]
\label{lem:appendix-half-square-admissibility}
If \((x,y,z)\in\mathcal A\) satisfies
\[
 x\ge\frac12,
 \qquad
 0\le z\le\frac12,
\]
then
\[
 y\le1-z.
\]
\end{lemma}

\begin{proof}
Suppose \(y>1-z\).  Then \(x+y>1\), so the selected supercritical cell of
\zcref{prop:new-exact-local-set} applies.  In the ordered half \(x\le y\)
its necessary polynomial is
\[
 F(x,y,z)=
 (x^2-1)z^2+(2xy^2+y)z+y^4-y^2\le0.
\]
It is nondecreasing in \(x\), its derivative in \(y\) is positive for
\(y\ge1-z\), and
\[
 F\left(\frac12,1-z,z\right)
 =\frac{z^2(2z-5)(2z-1)}4\ge0.
\]
This is a contradiction.  In the reflected half \(y<x\), apply the same
argument to \(F(y,x,z)\).  At the joining corner,
\[
 F(1-z,1-z,z)=z(1-2z)\ge0,
\]
so no selected feasible point is lost at equality.
\end{proof}

\begin{lemma}[Two-vertex replacement from scalar margins]
\label{lem:two-vertex-replacement}
Consider two adjacent V roles. Only the five pieces
\[
 e_{5,0},\quad e_{0,1},\quad e_{1,2},\quad r_0,\quad r_1
\]
may be affected by their replacement, and the shared edge is center-free.
Let $a,c,B,r\in[0,1)$ be the boundary reach on $e_{5,0}$, the own-radial
reach on $r_0$, the boundary reach on $e_{1,2}$, and the radial demand
on $r_1$ needed to overlap the unchanged center interval. If
\[
 a<\tfrac12,\qquad a+c<1,\qquad a+B<1,\qquad
 r<\max\{1-a,1-B\},
\]
there are two open unit equilateral replacements preserving these five
skeleton pieces, with actual boundary sums below one and no positive
adjacent support. No center-class formula is assumed.
\end{lemma}
\begin{proof}
Choose first $p_2\in(a,1-B)$ with $\max\{p_2,1-p_2\}>r$. This is
possible because the supremum on that interval is $\max\{1-a,1-B\}$.
Next choose $a<p_1<\min\{p_2,1/2,1-c\}$ and
\[
 0<\varepsilon<\min\{p_1-a,p_2-p_1,1-B-p_2,
                         1-p_1-c,\max(p_2,1-p_2)-r\}.
\]
All five margins are positive. Use the separate vertex charts
\[
 X_0(x,y)=V_0+x(V_5-V_0)+y(V_1-V_0),\qquad
 X_1(x,y)=V_1+x(V_0-V_1)+y(V_2-V_1).
\]
Both have metric $x^2+y^2-xy$. Define
\[
 \Delta_p^-=\conv\{(0,1-p),(1,1-p),(0,-p)\},\quad
 D^-_{p,\varepsilon}=\interior(\Delta_p^-)+(-\varepsilon,0)
 \quad(p\le1/2),
\]
\[
 \Delta_p^+=\conv\{(p,0),(p,1),(p-1,0)\},\quad
 D^+_{p,\varepsilon}=\interior(\Delta_p^+)+(0,-\varepsilon)
 \quad(p>1/2).
\]
Their edges have unit length. The minus half-planes are
$x>-\varepsilon$, $y<1-p$, $y>x+\varepsilon-p$; the plus half-planes
are $y>-\varepsilon$, $x<p$, $x>y+\varepsilon+p-1$.
Substitution on the axes and diagonal gives their actual reaches:
\[
\begin{array}{c|ccc}
 &A&B&C\\\hline
 D^-_{p,\varepsilon}&p-\varepsilon&1-p&1-p\\
 D^+_{p,\varepsilon}&p&1-p-\varepsilon&p
\end{array}
\]
Both contain the chart origin. On the adjacent support lines $(z,1)$ and
$(1,z)$ the half-planes are incompatible for $0\le z\le1$, so neither
has positive adjacent support.
Set $U'_0=X_0(D^-_{p_1,\varepsilon})$ and use
$U'_1=X_1(D^-_{p_2,\varepsilon})$ for $p_2\le1/2$, or
$U'_1=X_1(D^+_{p_2,\varepsilon})$ otherwise. The two charts and this
split cannot be identified. Both actual boundary sums are $1-\varepsilon$.

The first replacement reaches $p_1-\varepsilon>a$ on $e_{5,0}$ and
$1-p_1>c$ on $r_0$. On the shared edge the reaches sum to at least
\[
 (1-p_1)+(p_2-\varepsilon)>1,
\]
so the open traces overlap. The second reaches more than $B$ on
$e_{1,2}$ and $\max(p_2,1-p_2)>r$ on $r_1$. Thus all five pieces
remain covered; all other pieces are unchanged.
\end{proof}

\begin{proposition}[Vd1 input to the two-vertex replacement]
\label{prop:appendix-vd1-two-chart-replacement}
Assume a skeleton cover has a unique Vd1 role $T_0$ supplying $M_1$, a
unique supercritical role $T_1$, and all other V roles are nonsupercritical
Vd0. Suppose $e_{0,1}$ is center-free and no other V role has positive
trace on $r_0$ or $r_1$. The scalar replacement applies in CE1 and CE2
and yields a skeleton cover with $N'_+=0$.
\end{proposition}

\begin{proof}
Let \(A_0,B_0,C_0\) be the actual incoming, outgoing, and own-radial
reaches of the Vd1 role.  Its corner normal form and \zcref{lem:vd1-forward-orientation} supply
\(\vartheta\ge1\) and
\[
 \omega:=\sqrt{\vartheta^2+\vartheta+1}
\]
such that
\[
 C_0=\frac{\omega-A_0-\vartheta B_0}{\vartheta+1},
 \qquad
 \lambda=\frac{\vartheta(1-B_0)}{\vartheta+1},
\]
\[
 u_{\rm adj}
 =\frac{\omega-A_0-\vartheta B_0-1}{\vartheta}.
\]
The strict midpoint and one-sided-support conditions give
\[
 A_0+C_0<1,
 \qquad
 A_0<\lambda\le\frac12,
\]
\[
 B_0<\frac12,
 \qquad
 u_{\rm adj}<1-A_0.
\]
For the first inequality, evaluation on the closed midpoint face reduces
strictness to
\[
 (2\vartheta^2+4\vartheta+1)^2
 -4(\vartheta+1)^2(\vartheta^2+\vartheta+1)
 =4\vartheta^3+4\vartheta^2-4\vartheta-3>0.
\]
The remaining inequalities follow directly from the displayed endpoint
formulas and the strict midpoint inequalities.

The center-free shared edge and the supported interval give
\[
 A_1\ge1-B_0>\frac12,
 \qquad
 C_1\ge\lambda
\]
for the actual reaches of \(T_1\).  The selected supercritical component
has \(C_1\le1/2\), so
\zcref{lem:appendix-half-square-admissibility} yields
\[
 B_1\le1-C_1\le1-\lambda.
\]
Since \(A_0<\lambda\),
\[
 A_0+B_1<1.
\]

Define the C-triangle reach on \(r_1\), measured from \(O\), by
\[
 d_1^C:=\max\{x\in[0,1]:xV_1\in T_C\},
 \qquad
 c_1^{\rm req}:=1-d_1^C.
\]
Only \(U_C,U_1\), and the adjacent trace of \(U_0\) can contribute a
positive interval to \(r_1\).  Open skeleton coverage gives
\[
 c_1^{\rm req}<\max\{C_1,u_{\rm adj}\}
 \le\max\{1-B_1,1-A_0\}.
\]

Apply \zcref{lem:two-vertex-replacement} with
\[
 (a,c,B,r)=(A_0,C_0,B_1,c_1^{\rm req}).
\]
All four scalar conditions were proved above. The four untouched roles
remain nonsupercritical Vd0, while the two new actual boundary sums are
$1-\varepsilon<1$. The resulting skeleton cover contradicts N0.
The original C triangle is fixed; its class and the input gap rank were
not used in the construction. This operation does not claim to preserve
coverage of the interior of $H$.
\end{proof}
